%======================================================================
% فصل اول: مقدمه
%======================================================================
\chapter{مقدمه}
\label{ch:intro}

%----------------------------------------------------------------------
\section{کلیات}
\label{sec:intro-overview}
%----------------------------------------------------------------------

در دهه‌ی اخیر، \term{هوش مصنوعی}{Artificial Intelligence (AI)} تحولی
بنیادین در حوزه‌های مختلف فناوری اطلاعات ایجاد کرده است. ظهور
\term{مدل‌های زبانی بزرگ}{Large Language Models (LLM)} مانند
\lr{GPT-4o}، \lr{Gemini} و \lr{Claude}، امکانات بی‌سابقه‌ای را برای
توسعه‌دهندگان نرم‌افزار و سازمان‌ها فراهم کرده است. این مدل‌ها از طریق
\term{رابط‌های برنامه‌نویسی کاربردی}{Application Programming Interface (API)}
در دسترس عموم قرار گرفته‌اند و کاربردهای گسترده‌ای در حوزه‌هایی از جمله
دستیاران هوشمند، تولید محتوا، تحلیل داده، تبدیل گفتار به متن و تولید تصویر
یافته‌اند~\cite{openai2024}.

شرکت‌های بزرگ فناوری از جمله \lr{OpenAI}، \lr{Google} (با محصول
\lr{Gemini})، \lr{Anthropic} (با محصول \lr{Claude})، \lr{Groq} و
سرویس‌های واسط مانند \lr{OpenRouter} هر کدام رابط‌های برنامه‌نویسی
منحصربه‌فرد خود را ارائه می‌دهند. گرچه بسیاری از این رابط‌ها قابلیت‌های
مشابهی دارند، اما در ساختار درخواست، احراز هویت، فرمت پاسخ و سیستم
صورت‌حساب‌گذاری تفاوت‌های قابل توجهی دارند~\cite{groq2024, openrouter2024}.

این تنوع باعث شده که سازمان‌هایی که قصد دارند از چند سرویس هوش مصنوعی
به طور هم‌زمان بهره ببرند یا از یک ارائه‌دهنده به ارائه‌دهنده‌ی دیگری مهاجرت
کنند، با چالش‌های فنی و مدیریتی قابل توجهی روبه‌رو شوند. در این شرایط،
وجود یک \term{سامانه‌ی واسط}{Gateway} که بتواند ارتباط با سرویس‌های مختلف
را به صورت استاندارد مدیریت کند، نیاز ملموسی است.

این پروژه با هدف طراحی و پیاده‌سازی چنین سامانه‌ای تعریف شده است. سامانه‌ی
پیاده‌سازی‌شده به عنوان یک لایه‌ی میانجی بین
\term{سرویس‌گیرندگان}{Clients} و ارائه‌دهندگان مختلف هوش مصنوعی عمل
می‌کند و یک رابط برنامه‌نویسی استاندارد و یکپارچه را در اختیار
توسعه‌دهندگان قرار می‌دهد. علاوه بر این، قابلیت‌هایی مانند احراز هویت،
مدیریت کلیدهای رابط برنامه‌نویسی، کیف پول، خرید \term{اشتراک}{Subscription}،
صدور فاکتور و گزارش‌های مصرف نیز در این سامانه پیش‌بینی شده است.

%----------------------------------------------------------------------
\section{بیان مسئله}
\label{sec:problem}
%----------------------------------------------------------------------

در اکوسیستم کنونی هوش مصنوعی، هر ارائه‌دهنده‌ی سرویس رابط برنامه‌نویسی
مخصوص خود را تعریف می‌کند. برای مثال، \lr{OpenAI} یک ساختار مشخص برای
ارسال پیام و دریافت پاسخ دارد، \lr{Gemini} از \lr{API} بومی
\lr{Google} با فرمت متفاوتی استفاده می‌کند، و \lr{Groq} اگرچه سازگار با
فرمت \lr{OpenAI} است، اما \lr{API Key} و نقاط پایانی\footnote{\lr{Endpoint}}
جداگانه‌ای دارد. این تفاوت‌ها در چند سطح مشکل‌ساز هستند:

\begin{enumerate}
  \item \textbf{چندگانگی پیاده‌سازی:}
    توسعه‌دهنده‌ای که می‌خواهد از چند سرویس هوش مصنوعی استفاده کند، ملزم
    است برای هر کدام کد جداگانه‌ای بنویسد. این موضوع حجم کد را افزایش داده
    و نگهداری سیستم را دشوار می‌کند.

  \item \textbf{چالش مهاجرت:}
    در صورتی که سازمانی بخواهد از یک ارائه‌دهنده به ارائه‌دهنده‌ی دیگری
    مهاجرت کند—برای مثال از \lr{OpenAI} به \lr{Groq}—ناچار به تغییرات
    گسترده در کدپایه\footnote{\lr{Codebase}} خواهد بود.

  \item \textbf{پراکندگی مدیریت:}
    کلیدهای \lr{API}، \term{محدودیت نرخ}{Rate Limit}، گزارش مصرف و
    صورت‌حساب‌ها برای هر سرویس به صورت مجزا مدیریت می‌شوند. این پراکندگی
    دید جامعی از مصرف کل را از کاربر می‌گیرد.

  \item \textbf{کمبود ابزار مدیریت مالی:}
    اکثر ارائه‌دهندگان سرویس هوش مصنوعی امکانات محدودی برای مدیریت
    مصرف چندکاربره، کیف پول داخلی، صدور فاکتور یا خرید اشتراک دارند.
    این نیاز برای سازمان‌هایی که می‌خواهند دسترسی به هوش مصنوعی را
    برای کاربران داخلی خود مدیریت کنند، بیشتر احساس می‌شود.

  \item \textbf{نبود کنترل و نظارت متمرکز:}
    بدون یک سامانه‌ی واسط، کنترل دسترسی، نظارت بر مصرف، تعیین سقف
    مصرف و ثبت گزارش‌های تفصیلی برای هر ارائه‌دهنده باید به صورت جداگانه
    پیاده‌سازی شود.
\end{enumerate}

این مشکلات با گسترش استفاده از هوش مصنوعی در سازمان‌ها حادتر می‌شود.
راهکار پیشنهادی این پروژه، طراحی و پیاده‌سازی یک سامانه‌ی واسط است که
تمامی این چالش‌ها را از طریق یک لایه‌ی مرکزی و یکپارچه حل کند.

%----------------------------------------------------------------------
\section{انگیزه و ضرورت}
\label{sec:motivation}
%----------------------------------------------------------------------

انگیزه‌ی اصلی این پروژه از نیاز واقعی سازمان‌ها و توسعه‌دهندگانی نشأت
می‌گیرد که می‌خواهند از قابلیت‌های هوش مصنوعی بهره ببرند اما نمی‌خواهند
درگیر پیچیدگی‌های یکپارچه‌سازی چندین سرویس مجزا شوند. ضرورت چنین
سامانه‌ای را می‌توان از چند منظر بررسی کرد:

\textbf{از منظر فنی،} با داشتن یک رابط برنامه‌نویسی استاندارد و سازگار با
\lr{OpenAI}—که پرکاربردترین فرمت در این حوزه است—توسعه‌دهندگان می‌توانند
بدون تغییر در کد سمت مشتری، از ارائه‌دهندگان مختلف استفاده کنند. تنها با
تغییر فیلدهای \code{provider} و \code{model} در درخواست، سامانه
درخواست را به ارائه‌دهنده‌ی مناسب هدایت می‌کند.

\textbf{از منظر مالی،} یک سیستم متمرکز مدیریت کیف پول و صورت‌حساب،
امکان صدور فاکتور یکپارچه، ردیابی دقیق مصرف و کنترل هزینه‌ها را فراهم
می‌کند. این موضوع برای سازمان‌هایی که هزینه‌های هوش مصنوعی بخش قابل
توجهی از بودجه‌ی عملیاتی آن‌ها را تشکیل می‌دهد، از اهمیت ویژه‌ای
برخوردار است.

\textbf{از منظر امنیتی،} متمرکزسازی احراز هویت و مدیریت کلیدهای
\lr{API} در یک سامانه‌ی واحد، امکان اعمال سیاست‌های امنیتی منسجم—مانند
محدودیت نرخ درخواست، لیست سفید آدرس \lr{IP} و \term{ممیزی}{Audit}
عملیات حساس—را به صورت مرکزی فراهم می‌آورد.

\textbf{از منظر مقیاس‌پذیری،} الگوی خط لوله\footnote{\lr{Pipeline Pattern}}
به‌کاررفته در این سامانه، امکان افزودن مراحل جدید به پردازش درخواست—مثلاً
\term{مدار قطعه}{Circuit Breaker} یا \term{راهبری هوشمند}{Smart Routing}—را
بدون تغییر در ساختار اصلی فراهم می‌کند.

%----------------------------------------------------------------------
\section{اهداف پروژه}
\label{sec:objectives}
%----------------------------------------------------------------------

هدف اصلی این پروژه، طراحی و پیاده‌سازی یک سامانه‌ی واسط جامع برای
یکپارچه‌سازی سرویس‌های هوش مصنوعی است. این هدف کلی به اهداف عملیاتی
زیر تجزیه می‌شود:

\begin{enumerate}
  \item \textbf{ارائه‌ی رابط برنامه‌نویسی یکپارچه:}
    پیاده‌سازی مجموعه‌ای از \term{نقاط پایانی}{Endpoints} سازگار با
    استاندارد \lr{OpenAI}، شامل: تکمیل گفتگو\footnote{\lr{Chat Completions}}،
    \term{جاسازی متن}{Embeddings}، تولید تصویر\footnote{\lr{Image Generation}}،
    تبدیل گفتار به متن\footnote{\lr{Audio Transcription}} و
    تبدیل متن به گفتار\footnote{\lr{Text-to-Speech}}.

  \item \textbf{پشتیبانی از چند ارائه‌دهنده:}
    پیاده‌سازی \term{آداپتور}{Adapter} برای ارائه‌دهندگان \lr{OpenAI}،
    \lr{Gemini}، \lr{Groq} و \lr{OpenRouter} با امکان افزودن ارائه‌دهندگان
    جدید به صورت \term{اتصال‌پذیر}{Pluggable}.

  \item \textbf{احراز هویت و مدیریت دسترسی:}
    پیاده‌سازی جریان احراز هویت مبتنی بر رمز یک‌بارمصرف برای ورود به
    داشبورد مدیریتی، و مدیریت کلیدهای \lr{API} برای درخواست‌های برنامه‌ای.

  \item \textbf{سیستم صورت‌حساب:}
    پیاده‌سازی کیف پول کاربری، ثبت دقیق مصرف توکن و منابع غیرتوکنی،
    صدور فاکتور و مدیریت اشتراک‌های ماهانه.

  \item \textbf{گزارش‌دهی و مانیتورینگ:}
    فراهم کردن گزارش‌های تفصیلی مصرف، دفترکل تراکنش‌های کیف پول و
    لیست فاکتورها، همراه با ارسال گزارش‌های خط لوله به سرویس‌های
    مدیریت گزارش خارجی.

  \item \textbf{آزمون و مستندسازی:}
    نوشتن مجموعه‌ی آزمون‌های واحد و ویژگی برای تمامی ماژول‌ها، و
    تولید مستندات \lr{OpenAPI} با استفاده از ابزار \lr{Swagger}.
\end{enumerate}

%----------------------------------------------------------------------
\section{دستاوردهای پروژه}
\label{sec:contributions}
%----------------------------------------------------------------------

پروژه‌ی حاضر دستاوردهای زیر را به عنوان خروجی اصلی ارائه می‌دهد:

\begin{itemize}
  \item \textbf{سامانه‌ی واسط عملیاتی:} یک نرم‌افزار سمت سرور کاملاً
    عملیاتی که درخواست‌های هوش مصنوعی را از سرویس‌گیرندگان دریافت و به
    چهار ارائه‌دهنده‌ی \lr{OpenAI}، \lr{Gemini}، \lr{Groq} و
    \lr{OpenRouter} هدایت می‌کند.

  \item \textbf{رابط برنامه‌نویسی سازگار با \lr{OpenAI}:}
    پیاده‌سازی پنج نقطه‌ی پایانی هوش مصنوعی با فرمت سازگار با استاندارد
    \lr{OpenAI}، به گونه‌ای که سرویس‌گیرندگان موجود با حداقل تغییر
    می‌توانند به این سامانه متصل شوند.

  \item \textbf{معماری مبتنی بر خط لوله:}
    پیاده‌سازی چهارده مرحله‌ی پردازشی در قالب الگوی خط لوله‌ی \lr{Laravel}،
    که هر مرحله مسئولیتی مشخص و مستقل دارد.

  \item \textbf{سیستم صورت‌حساب دقیق:}
    اندازه‌گیری مصرف توکن با استفاده از توکن‌ساز واقعی \lr{BPE}، پشتیبانی
    از صورت‌حساب‌گذاری منابع غیرتوکنی (تصویر، صوت) و مدیریت کیف پول با
    تراکنش‌های ایمن.

  \item \textbf{داشبورد \lr{Playground}:}
    یک رابط کاربری \lr{Vue.js} برای آزمایش تعاملی نقاط پایانی سامانه، با
    امکان انتخاب ارائه‌دهنده، مدل، کلید \lr{API} و مشاهده‌ی پاسخ.

  \item \textbf{پوشش آزمون:}
    مجموعه‌ای از آزمون‌های واحد و ویژگی که جریان‌های اصلی سامانه را
    پوشش می‌دهند.
\end{itemize}

%----------------------------------------------------------------------
\section{تغییرات نسبت به پروپوزال}
\label{sec:changes}
%----------------------------------------------------------------------

در \term{پروپوزال}{Proposal} اولیه‌ی این پروژه، بسته‌ی فناوری پیش‌بینی‌شده
شامل زبان \lr{Python}، چارچوب‌های \lr{FastAPI} یا \lr{Django}، و پایگاه داده‌ی
\lr{PostgreSQL} بود. در مرحله‌ی اجرا، این بسته‌ی فناوری به دلایل فنی و
عملی تغییر یافت و پیاده‌سازی نهایی با فناوری‌های زیر انجام شد:

\begin{itemize}
  \item \textbf{چارچوب \lr{Laravel} (زبان \lr{PHP})} — به عنوان چارچوب
    اصلی \term{سمت سرور}{Server-side}.
  \item \textbf{پایگاه داده‌ی \lr{SQL Server}} — به عنوان منبع اصلی ذخیره‌ی
    داده.
  \item \textbf{\lr{Redis}} — به عنوان \term{موتور کش}{Cache Engine} و
    مدیریت صف‌های پردازش پس‌زمینه.
  \item \textbf{\lr{Vue.js}} — به عنوان چارچوب رابط کاربری برای
    \term{میدان بازی}{Playground}.
\end{itemize}

اهداف اصلی پروژه—شامل یکپارچه‌سازی چند ارائه‌دهنده‌ی هوش مصنوعی، سیستم
احراز هویت، مدیریت کلیدهای \lr{API}، کیف پول، اشتراک و گزارش‌دهی—بدون
تغییر باقی ماندند و کاملاً پیاده‌سازی شدند.

%----------------------------------------------------------------------
\section{ساختار پایان‌نامه}
\label{sec:structure}
%----------------------------------------------------------------------

مابقی این پایان‌نامه به صورت زیر سازماندهی شده است:

\begin{description}
  \item[\textbf{فصل دوم — مروری بر پیشینه:}]
    ابتدا مفاهیم پایه‌ی سامانه‌های واسط و رابط‌های برنامه‌نویسی هوش مصنوعی
    بررسی می‌شوند. سپس سیستم‌های مشابه از جمله \lr{Groq}، \lr{OpenRouter}،
    \lr{Anyscale Endpoints} و \lr{OpenAI} تحلیل و مقایسه می‌شوند.

  \item[\textbf{فصل سوم — طراحی سیستم:}]
    معماری کلی سامانه، مدل داده، طراحی \lr{API}، الگوی خط لوله، استراتژی کش
    و طراحی سیستم صورت‌حساب تشریح می‌شوند.

  \item[\textbf{فصل چهارم — پیاده‌سازی:}]
    جزئیات پیاده‌سازی هر ماژول از جمله خط لوله‌ی سامانه‌ی واسط، آداپتورهای
    ارائه‌دهندگان، احراز هویت، مدیریت کلیدهای \lr{API} و سیستم صورت‌حساب
    ارائه می‌شود.

  \item[\textbf{فصل پنجم — ارزیابی:}]
    نتایج آزمون‌های واحد و ویژگی، ارزیابی عملکرد و مقایسه با اهداف اولیه
    ارائه می‌شوند.

  \item[\textbf{فصل ششم — نتیجه‌گیری:}]
    دستاوردهای کلی پروژه جمع‌بندی شده و پیشنهادهایی برای کارهای آتی ارائه
    می‌شود.
\end{description}
